\section{Related Work}
\label{sec:related_work}

\subsection{2D Visual Decoding from Brain Activity}


Visual decoding from brain activity \cite{lahner2024modeling,yi2024learning,chen2023seeing,chen2024eegformer,chen2024cinematic,zhang2022neural} has gained substantial attention in computer vision and neuroscience, emerging as an effective technique for understanding and analyzing human visual perception mechanisms.
Early approaches predominantly utilized Convolutional Neural Networks (CNNs) and Generative Adversarial Networks (GANs) to model brain activity signals and interpret visual information~\cite{horikawa2017generic,shen2019deep,beliy2019voxels,lin2022mind,ozcelik2022reconstruction,huang2021fmri}. 
Recently, the utilization of newly-emerged diffusion models \cite{ho2020denoising,nichol2021improved} and vision-language models \cite{li2023blip,zhu2023minigpt} has advanced visual generation from various neural signals including fMRI \cite{chen2023seeing,takagi2023high,scotti2024reconstructing,sun2024contrast,scotti2024mindeye2}
and EEG \cite{song2023decoding,bai2023dreamdiffusion,li2024visual,singh2023eeg2image}. 
These methods typically perform contrastive alignment \cite{wang2020understanding} between 
neural signal embeddings and image or text features derived from the pre-trained CLIP model~\cite{radford2021learning}. Subsequently, the aligned neural embeddings are sent into diffusion model to conditionally reconstruct images that correspond to the visually-evoked brain activity. 
Apart from static images,
research has begun to extend these approaches to the reconstruction of video information from fMRI data, further advancing the field \cite{chen2024cinematic,sun2024neurocine,wen2018neural,wang2022reconstructing,kupershmidt2022penny}.
Though impressive, these methods, limited to 2D visual perception, fall short of capturing the full depth of human 3D perceptual experience in real-world environments. Our method attempts to expand the scope of brain visual decoding to three dimensions by reconstructing 3D objects from real-time EEG signals.

\subsection{3D Reconstruction from fMRI}

Reconstructing 3D objects from brain signals holds significant potential for advancing both brain analysis applications and our understanding of the brain’s visual system. To achieve this goal, several works~\cite{gao2023mind,fMRI-3D} have made initial strides in 3D object reconstruction from fMRI, yielding promising results in interpreting 3D spatial structures. 
Mind-3D~\cite{gao2023mind} proposes the first dataset of paired fMRI and 3D shape data and develops a diffusion-based framework to decode 3D shape from fMRI signals. A subsequent work, fMRI-3D~\cite{fMRI-3D}, expanded the dataset to include a broader range of categories across five subjects. 

However, there are several limitations in previous task setup, which prevent it from simulating real-time, natural 3D perception scenarios. First, fMRI equipment is unportable, expensive, and difficult to operate, potentially hindering its application in BCIs (Brain-Computer-Interface). Besides its high acquisition cost, fMRI is limited by its inherently low temporal resolution, which hinders real-time responsiveness to dynamic stimuli. Second, existing brain 3D reconstruction methods focus exclusively on reconstructing 3D shape of objects, neglecting color and appearance information that are crucial in real-world perception.
To address these challenges, 
we introduce a 3D visual decoding framework based on EEG signals, along with 
a new dataset for paired EEG and colored objects. 
To the best of our knowledge, this is the first work to interpret 3D objects from EEG signals, offering a comprehensive dataset, benchmarks, and decoding framework.


\vspace{-1mm}
\subsection{Diffusion Models}



Diffusion model has recently emerged as a powerful generative framework known for high-quality image synthesis capabilities. Inspired by non-equilibrium thermodynamics, the diffusion models are formulated as Markov chains. The model first progressively corrupts the target data distribution by adding noise until it conforms to a standard Gaussian distribution, and subsequently generates samples by predicting and reversing the noise process through network learning \cite{ho2020denoising,nichol2021improved}. The diffusion model, along with its variants, has been extensively applied to tasks such as image generation \cite{rombach2022high,kim2022diffusionclip,saharia2022photorealistic,ramesh2021zero} and image editing \cite{yu2024accelerating,zhang2023adding}.


Building on 2D image generation, \cite{luo2021diffusion} and \cite{zhou20213d} extended pixel-based approaches to 3D coordinates, enabling the generation of point clouds. This has spurred further research into 3D generation \cite{ren2024tiger,vahdat2022lion}, text-to-3D reconstruction \cite{poole2022dreamfusion,nichol2022point}, and 2D-to-3D generation \cite{melas2023pc2,wu2023sketch,long2024wonder3d}, demonstrating their capability to capture intricate spatial structures and textures of 3D objects.
In our study, we extend the 3D diffusion model to brain activity analysis, reconstructing colored 3D objects from EEG signals.